\documentclass[11pt,a4paper]{article}

\usepackage[margin=1in]{geometry}
\usepackage{amsmath,amssymb}
\usepackage{graphicx}
\usepackage{booktabs}
\usepackage{enumitem}
\usepackage{hyperref}
\usepackage{xcolor}
\usepackage{float}
\usepackage{microtype}

\setlength{\parindent}{0pt}
\setlength{\parskip}{0.5em}
\setlist[itemize]{nosep,leftmargin=*}
\setlist[enumerate]{nosep,leftmargin=*}

\title{\textbf{ID6002W : RL Course Project Report}\\[0.4em]
\large Industrial Inventory Control}
\author{Student Name \quad | \quad Roll Number}
\date{}

\begin{document}

\maketitle

\begin{center}
\fbox{\parbox{0.92\textwidth}{
\textbf{Recommended report length: 4--5 pages including references.}
Keep the report concise and evidence-driven. Focus on what you tried, why you tried it, what worked, what did not work, and what you learned from the experiments.
}}
\end{center}

\section{Problem and Assigned Configuration}

Briefly describe the inventory-control problem in your own words. Do not repeat the full project statement.

\begin{itemize}
    \item \textbf{Variant ID:} \texttt{<enter variant ID>}
    \item \textbf{Config fingerprint:} \texttt{<enter fingerprint>}
    \item \textbf{Demand multiplier profile:} \texttt{<enter values>}
    \item \textbf{Initial inventory profile:} \texttt{<enter values>}
    \item \textbf{Lead-time delay profile:} \texttt{<enter values>}
\end{itemize}

\section{Techniques Attempted}

List each distinct technique that you tried. For each technique, summarize the state/observation variables used, action representation, learning method, and any important design choices.

\begin{table}[H]
\centering
\small
\begin{tabular}{p{0.18\textwidth}p{0.3\textwidth}p{0.5\textwidth}}
\toprule
\textbf{Technique} & \textbf{Observations Used} & \textbf{Key Design Choices} \\
\midrule
Technique 1 &  &    \\
Technique 2 &  &    \\
Technique 3 &  &    \\
\bottomrule
\end{tabular}
\caption{Summary of techniques attempted. Add or remove rows as required.}
\end{table}

\section{Training and Experimental Setup}

For the five top techniques, briefly report the experimental setup.

\begin{itemize}
    \item Number of training episodes / steps
    \item Discount factor, learning rate, and other important hyperparameters
    \item Exploration strategy, if applicable
    \item Reward shaping, if used
    \item Random seeds used for controlled comparisons
\end{itemize}

Do not list every hyperparameter unless it materially affects the result.

\section{Results}

Report the results of the top 5 techniques. Use the same evaluation protocol when comparing techniques.
\begin{table}[H]
\centering
\small
\begin{tabular}{lrrrr}
\toprule
\textbf{Technique} & \textbf{Mean Cost} & \textbf{Best Cost} & \textbf{Notes} \\
\midrule
Technique 1 &  &  &  \\
Technique 2 &  &  &  \\
Technique 3 &  &  &  \\
\bottomrule
\end{tabular}
\caption{Validation performance before leaderboard submission.}
\end{table}

Include at most two plots that materially help explain the results. Examples include training return versus episode, validation cost across techniques, or cost breakdown across scenarios.

% Example figure block:
% \begin{figure}[H]
%     \centering
%     \includegraphics[width=0.75\textwidth]{training_curve.png}
%     \caption{Training performance of the selected technique.}
% \end{figure}

\section{Analysis and Observations}

Discuss the most important findings from your experiments. Focus on evidence rather than general descriptions of the algorithms.

Suggested points:
\begin{itemize}
    \item Which technique performed best and why?
    \item Which observations were most useful?
    \item Did demand history or the arrival pipeline materially improve the policy?
    \item How sensitive was performance to demand variation or lead-time delays?
    \item What failure modes did you observe, such as stockouts, excessive holding, or discards?
\end{itemize}

\section{Key Learning}

Summarize the 3--5 most important technical lessons from the project.

\begin{enumerate}
    \item 
    \item 
    \item 
\end{enumerate}

\section*{References}

List only references that were actually used, such as papers, documentation, libraries, repositories, or other technical sources.


\end{document}
